\section{Conclusion}

We studied lightweight soft-prompt routing as a physical model-selection
mechanism for semantic operators. \system contributes two methodological
components: disagreement-guided acquisition that concentrates a limited human
gold-label budget on routing-critical candidates after a full paired screen,
and a cost-derived score whose inverse expected route cost supplies the WCE
weight. The acquired outcomes train a frozen-encoder soft-prompt router with
only 3,074--13,826 task-specific trainable parameters and no autoregressive
routing step. Reflective discrete optimization is a separate, substantially
larger alternative used to test whether the supervision transfers across
router classes.

On the Computer Science pool, labeling half of the screened candidates after
disagreement prioritization recovers every Mini-only and Nano-only case, whereas
random selection misses nearly half. This result supports the acquisition
mechanism while making its cost boundary explicit: it saves human annotations,
not Nano or Mini screening calls. The final matrix will verify the combined
disagreement-guided and inverse-expected-cost WCE pipeline under the shared cost
score.
The locked fresh set supplies the central warning: small validation splits and
rare critical cases make apparently safe savings unstable. Neither additional
soft-prompt capacity nor a more elaborate discrete policy fixes calibration by
itself.

The broader systems lesson is that an LLM router should be evaluated like a
learned query optimization component, not like a standalone classifier. Its
statistics must be outcome-specific, its operating point must be calibrated
without test reuse, its overhead must enter the cost model, and its quality must
be confirmed on data that was not part of prompt development. Completing the
pre-registered multi-operator evaluation will determine whether \system becomes
a deployable optimizer or remains a cautionary benchmark for prompt-based
routing.
